% !TEX root = ../monograph.tex
\chapter{Lattice gauge theory and the continuum problem}
\label{ch01:lattice}

The starting point is a probability measure on finitely many group-valued
variables. Its definition requires neither gauge fixing nor a continuum
functional integral. The construction problem is to choose its coupling and
its local observables as the lattice is refined, and to control their
correlations uniformly enough to remove the cutoff.

Two limits have different roles. At a fixed physical volume, the lattice
spacing will tend to zero along a prescribed sequence. Only afterwards will
the volume tend to infinity, along a subsequence. Keeping these operations
separate is essential: compactness can produce a subsequence, but cannot by
itself give convergence along every ultraviolet cutoff.

\section{The gauge group and its normalization}
\label{ch01:gauge}

Fix a closed subgroup $G\subset U(n_G)$ in a specified faithful unitary
representation. Its real Lie algebra $\gLie$ consists of skew-Hermitian
matrices, and we assume that $\gLie$ is simple and nonabelian. The group itself
need not be connected. We use the invariant inner product
\begin{equation}
 \langle X,Y\rangle_{\gLie}=-\frac1{n_G}\operatorname{Re}\tr(XY),
 \qquad X,Y\in\gLie.
 \label{ch01:metric}
\end{equation}
The choice of representation is part of the normalization of the model,
not an inessential choice of matrix coordinates.

Write $d_{\gLie}=\dim_{\RR}\gLie$ and fix once and for all an arbitrary orthonormal basis $(e_\alpha)$
for \eqref{ch01:metric}. The adjoint Casimir is the positive number
\begin{equation}
 C_{\gLie}=-\frac1{d_{\gLie}}\sum_{\alpha=1}^{d_{\gLie}}
       \tr_{\gLie}\bigl(\operatorname{ad}(e_\alpha)^2\bigr).
 \label{ch01:casimir}
\end{equation}
Here $\operatorname{ad}(X)Y=[X,Y]$ and $\tr_{\gLie}$ is the real operator
trace. Invariance of the inner product makes $\operatorname{ad}(X)$
skew-adjoint. Simplicity implies the familiar identity
$-\tr_{\gLie}(\operatorname{ad}(X)\operatorname{ad}(Y))
=C_{\gLie}\langle X,Y\rangle_{\gLie}$; nonabelianity gives $C_{\gLie}>0$.
We shall use the constants
\begin{equation}
 b_0=\frac{11C_{\gLie}}{48\pi^2},\qquad
 \kappa_G=2b_0=\frac{11C_{\gLie}}{24\pi^2}.
 \label{ch01:slope}
\end{equation}
Their role in the construction will be to describe the leading logarithmic
growth of the inverse coupling. Defining these constants does not yet
establish that growth for the nonperturbative lattice model.

For $G=SU(n)$ in its defining representation, our metric is
$-\operatorname{Re}\tr(XY)/n$, and
\begin{equation}
 d_{\gLie}=n^2-1,\qquad C_{\gLie}=2n^2,\qquad
 \kappa_G=\frac{11n^2}{12\pi^2}.
 \label{ch01:su-n}
\end{equation}
The conventional Wilson parameter is $\beta=2n/g^2$. Thus
$\beta\sim\kappa_G\log(1/a)$ corresponds to
$g^{-2}\sim(11n/24\pi^2)\log(1/a)$.
This is the normalization of the one-loop coefficient of
Gross--Wilczek and Politzer~\cite{GrossWilczek1973,Politzer1973}.
The exact control of the error in that law is a later construction problem.

\section{Links, holonomies and the Wilson measure}
\label{ch01:wilson-section}

For an integer $s\ge1$, set
\[
 \Lambda_s=(\ZZ/s\ZZ)^4,\qquad
 E_s=\Lambda_s\times\{1,2,3,4\}.
\]
The element $(y,\mu)$ of $E_s$ represents the oriented link from $y$ to
$y+e_\mu$. A configuration is $U\in G^{E_s}$. The variable on the reversed
link is, by definition, $U(y,\mu)^{-1}$. A path holonomy is the ordered
product of its link variables, with this inverse convention on backward
steps. We use row transport: multiplying a row vector at the initial point
by that ordered product transports it to the final point.

The elementary positively oriented plaquette in the $(\mu,\nu)$-plane,
with $\mu<\nu$, has holonomy
\begin{equation}
 U_{y;\mu\nu}=U(y,\mu)U(y+e_\mu,\nu)
                 U(y+e_\nu,\mu)^{-1}U(y,\nu)^{-1}.
 \label{ch01:plaquette}
\end{equation}
Define the plaquette energy and the normalized Hilbert--Schmidt distance from the identity by
\begin{equation}
 w(u)=1-\frac1{n_G}\operatorname{Re}\tr u,
 \qquad
 \chi(u)=\frac{\|u-1\|_{\mathrm F}}{\sqrt{n_G}},
 \qquad \chi(u)^2=2w(u).
 \label{ch01:chord}
\end{equation}
Both functions are invariant under conjugation and inversion; they vanish
at the identity. The last identity follows directly from $u^*u=1$.

\begin{definition}[Wilson measure]
\label{ch01:wilson}
Let $dh$ be probability Haar measure on $G$. The Wilson energy and measure
on $\Lambda_s$ are
\begin{align}
 E_s^{\mathrm W}(U)&=\sum_{y\in\Lambda_s}\sum_{\mu<\nu}w(U_{y;\mu\nu}),
 \label{ch01:wilson-energy}\\
 d\mu_{s,\beta}(U)&=Z_{s,\beta}^{-1}e^{-\beta E_s^{\mathrm W}(U)}
                     \prod_{e\in E_s}dh(U_e),
 &Z_{s,\beta}&=\int e^{-\beta E_s^{\mathrm W}(U)}\prod_{e\in E_s}dh(U_e).
 \label{ch01:wilson-law}
\end{align}
\end{definition}
Every unoriented coordinate square is counted once, through $\mu<\nu$.
There is no additional factor of $1/2$ in \eqref{ch01:wilson-energy}.
The configuration space is compact and the integrand defining $Z_{s,\beta}$
is continuous and strictly positive, so $0<Z_{s,\beta}<\infty$ for every
real $\beta$. The reflection-positivity arguments below require
$\beta\ge0$.

A site gauge transformation $g:\Lambda_s\to G$ acts by
\begin{equation}
 U^g(y,\mu)=g(y)U(y,\mu)g(y+e_\mu)^{-1}.
 \label{ch01:gauge-action}
\end{equation}
If $\gamma$ is a path from $y$ to $z$, its holonomy transforms as
$U_\gamma^g=g(y)U_\gamma g(z)^{-1}$. Thus a closed holonomy transforms by
conjugation. Together with left and right invariance of Haar measure, this
proves gauge invariance of the energy and of the law. These are the basic
properties of Wilson's lattice formulation~\cite{Wilson1974}; no gauge
fixing is involved in defining the probability measure.

\begin{definition}[Locality]
\label{ch01:locality}
A function of $U$ is local on $X\subset\Lambda_s$ if it depends only on
links whose two endpoints belong to $X$. It is gauge invariant if its value
is unchanged by \eqref{ch01:gauge-action} for every site gauge transformation.
A finite stencil is a finite subset of $\ZZ^4$, placed in the torus by
translation followed by reduction modulo $s$.
\end{definition}
Locality and gauge invariance are different requirements. A single link
matrix entry is local but usually not gauge invariant. A traced holonomy
around a fixed lattice loop satisfies both.

Gauge transformations also explain why small fields must be defined
through plaquettes rather than individual links. The pure-gauge
configuration $U(y,\mu)=g(y)g(y+e_\mu)^{-1}$ can have links far from the
identity while every plaquette holonomy is the identity. Later,
small-field regions will be specified by bounds on the gauge-invariant
quantity $\chi(U_{y;\mu\nu})$ from \eqref{ch01:chord}.

The connection with classical curvature is most cleanly stated using exact
parallel transport. For a smooth $\gLie$-valued connection $A$, let the
transport on a segment from $x$ to $x+ae_\mu$ solve
\[
 u'(t)=u(t)aA_\mu(x+tae_\mu),\qquad u(0)=1,\qquad 0\le t\le1.
\]
The plaquette product then has the expansion
\begin{equation}
 U_{x;\mu\nu}=1+a^2F_{\mu\nu}(x)+O(a^3),
 \quad
 F_{\mu\nu}=\partial_\mu A_\nu-\partial_\nu A_\mu+[A_\mu,A_\nu].
 \label{ch01:classical-curvature}
\end{equation}
Consequently
$w(U_{x;\mu\nu})=\frac12a^4|F_{\mu\nu}(x)|_{\gLie}^2+O(a^5)$.
This familiar classical expansion motivates the action. It is not a
regularity assertion about typical lattice configurations, nor about an
eventual continuum random connection.

\section{Physical scales and lattice symmetries}
\label{ch01:scales-section}

We keep distances in lattice units distinct from physical distances.
For $y\in\Lambda_s$, let $r_s(y)$ be its unique representative in
$(-s/2,s/2]^4\cap\ZZ^4$. The periodic distance in lattice units is
\[
 \operatorname{dist}_s(y,z)=
       \min_{k\in s\ZZ^4}|\widetilde y-\widetilde z+k|_\infty,
\]
where any lifts $\widetilde y,\widetilde z$ may be used. This distance is
translation invariant and invariant under signed coordinate permutations.
At a centered-cell boundary the chosen representative itself need not
transform as an ordinary vector; periodic statements are unaffected.

Let $H_4$ denote the finite group of signed permutation matrices. An affine
lattice isometry $y\mapsto hy+z$, $h\in H_4$, acts on configurations by
pulling back oriented paths, using inverse links when orientation is
reversed. Plaquette traces and product Haar measure are preserved, so the
Wilson law is invariant under these isometries. Full Euclidean rotation
invariance is not a lattice symmetry. Its recovery in the continuum is a
separate part of the argument.

\begin{definition}[Scale geometry]
\label{ch01:scales}
Choose odd integers $L_m\ge3$ and set
\begin{equation}
 D_0=1,\qquad D_{m+1}=L_mD_m,
 \qquad a_m=D_m^{-1},\qquad \tau_m=\log D_m.
 \label{ch01:scale-sequence}
\end{equation}
Choose strictly increasing physical periods $R_p$, each an integer multiple
of four. At scale $m$ and volume index $p$, the lattice has
$s_{m,p}=R_pD_m$ sites per side and physical spacing $a_m$.
\end{definition}
Its physical side is exactly $a_ms_{m,p}=R_p$. Increasing $m$ means refining
the lattice; integrating out one scale goes from $m+1$ to $m$. The index
$N$ will denote the finest scale in a finite renormalization trajectory.
At cutoff $N$, a lattice site has physical placement
\begin{equation}
 \widehat y=a_N r_{R_pD_N}(y)\in(-R_p/2,R_p/2]^4.
 \label{ch01:placement}
\end{equation}
Further growth and divisibility requirements on these sequences will be
imposed when the estimates that require them have been established.

Odd block ratios give centered blocks without an asymmetric choice of
boundary sites. At a step of ratio $L$, every fine site can be written
uniquely as $Ly+z$, with a coarse site $y$ and
$z\in\{-(L-1)/2,\ldots,(L-1)/2\}^4$. The resulting parent map commutes
with signed coordinate permutations. The corresponding coarse point and
its fine representative have the same physical position.

If an observable uses a fixed stencil $\Sigma\subset\ZZ^4$, its physical
support has diameter at most $a_N\operatorname{diam}(\Sigma)$. Such an
observable becomes local at a point as $N\to\infty$. This is different
from a loop of fixed physical size, whose number of links grows with
$N$. The latter is not included merely by proving estimates for every
fixed stencil.

\section{Reflection positivity}
\label{ch01:reflection}

Reflection positivity is the link between Euclidean correlations and a
positive Hilbert-space inner product. On the lattice it also gives
chessboard estimates, which convert a bound on many reflected copies of a
local observable into a bound on a correlation of separated observables.
We use these classical results as inputs~\cite{OsterwalderSeiler1978,
FrohlichIsraelLiebSimon1978,Seiler1982}.

Here are the support and orientation conventions for the finite-volume
positivity statement. Fix an even side $s=2h$ and one coordinate, temporarily
called $t$. A site reflection acts by $t\mapsto-t$ modulo $2h$; a link
reflection acts by $t\mapsto-1-t$. Their positive sectors are respectively
$1\le t<h$ and $0\le t<h$. In the first case the two fixed time slices
are excluded. In either case observables in the positive algebra depend
only on links whose two endpoints are in the indicated sector.

If the reflected coordinate is $\mu$ and the reflection of sites is
$\vartheta$, the reflected configuration is
\begin{equation}
 (\Theta U)_\nu(y)=
 \begin{cases}
 U_\nu(\vartheta y),&\nu\ne\mu,\\
 U_\mu(\vartheta y-e_\mu)^{-1},&\nu=\mu.
 \end{cases}
 \label{ch01:reflection-links}
\end{equation}
The inverse in the second line is essential. For a complex observable set
$(\theta F)(U)=\overline{F(\Theta U)}$.

\begin{proposition}[Finite-volume reflection positivity]
\label{ch01:lattice-rp}
Let $\beta\ge0$. For either of the two reflection cuts above and every
bounded Borel gauge-invariant observable $F$ in the positive algebra,
\begin{equation}
 \int (\theta F)(U)F(U)\,d\mu_{2h,\beta}(U)\ge0.
 \label{ch01:rp}
\end{equation}
The same conclusion holds after translating the cut or choosing another
coordinate direction.
\end{proposition}
\begin{proof}
This is the finite-torus form of lattice gauge reflection positivity
\cite{OsterwalderSeiler1978,Seiler1982}. Polarization gives a positive
semidefinite pairing between finite linear combinations of positive-side
observables. The support restriction is part of the statement: it is not
an assertion about arbitrary dependence on links crossing the reflection
planes. Positivity of the Boltzmann density, by itself, would not imply
\eqref{ch01:rp}.
\end{proof}

For clarity, one version of the chessboard construction can be described
without observables on seams. Write $s=2q\ell$ and cut each coordinate into
$2q$ intervals of length $\ell$. Let the base cell consist of sites with
all coordinates strictly between $0$ and $\ell$. An observable local on
this cell reads no seam links. Reflecting the cell repeatedly through the
planes at multiples of $\ell$ gives its copies in all
$M=(2q)^4$ cells. Denote the resulting real pullback of a function $F$ to
cell $C$ by $\mathcal R_CF$; the orientation of links is always transformed
as in \eqref{ch01:reflection-links}. For a bounded gauge-invariant $F$ on
the base cell, put
\begin{equation}
 \|F\|_{\mathrm{cb},\mathscr T}
 =\left[\int\prod_{C\in\mathscr T}\mathcal R_C|F|(U)
                 \,d\mu_{s,\beta}(U)\right]^{1/M},
 \label{ch01:chessboard-root}
\end{equation}
where $\mathscr T$ is this specified tiling. The quantity in brackets is
nonnegative. This notation records the tiling as well as the observable;
it is not a volume-independent norm by definition.

The chessboard estimate bounds, for one base-cell observable $F_C$ per cell,
\begin{equation}
 \left|\int\prod_{C\in\mathscr T}\mathcal R_CF_C\,d\mu_{s,\beta}\right|
 \le\prod_{C\in\mathscr T}\|F_C\|_{\mathrm{cb},\mathscr T}.
 \label{ch01:chessboard}
\end{equation}
Apply the abstract chessboard theorem
\cite[Theorem~4.1]{FrohlichIsraelLiebSimon1978} to the nonnegative
observables $|F_C|$, one coordinate at a time. The slots are the $2q$
strips in that coordinate: translation by $\ell$ gives cyclicity, and
Proposition~\ref{ch01:lattice-rp} gives its reflection Cauchy--Schwarz
hypothesis across the strip seams. The involution reflects the local
strip coordinates, including link orientation. Repeating in the other
three coordinates gives the reflected copies in
\eqref{ch01:chessboard-root}. No power-of-two restriction on $2q$ is
needed. Empty cells are filled with the constant one.
More flexible placements used later require explicit control of support,
seams and divisibility. They are not consequences of changing the tiling
symbol in \eqref{ch01:chessboard-root} without checking those conditions.

\section{Correlations and the order of the limits}
\label{ch01:limits}

Recall that $a_N$ is the lattice spacing and $R_p$ the physical side of the
torus, as in Definition~\ref{ch01:scales}. For now let $\mathcal T$ be an
arbitrary set of local-observable labels. A family of insertions assigns to
$v\in\mathcal T$ and $y\in\Lambda_{R_pD_N}$ a continuous gauge-invariant
function $I_{N,p}(v,y;U)$. For each fixed $v$, require locality on a translate
of a finite stencil $\Sigma_v$ independent of $N$ and $p$.
The next chapter constructs such local representatives; their necessary
cutoff-dependent normalization is a later issue.

Choose bare couplings $\beta_{N,p}$ and abbreviate
$\mu_{N,p}=\mu_{R_pD_N,\beta_{N,p}}$, the Wilson probability from
\eqref{ch01:wilson-law}. For a list $v=(v_1,\ldots,v_n)$, define
\begin{align}
 M^{\mathrm{lat}}_{N,p}(v;y)
   &=\int\prod_{i=1}^n I_{N,p}(v_i,y_i;U)\,d\mu_{N,p}(U),\\
 C^{\mathrm{lat}}_{N,p}(v;f)
   &=a_N^{4n}\sum_{y\in\Lambda_{R_pD_N}^n}
                   f(\widehat y_1,\ldots,\widehat y_n)
                   M^{\mathrm{lat}}_{N,p}(v;y).
 \label{ch01:lattice-smear}
\end{align}
The factor $a_N^{4n}$ is the physical volume element for $n$ insertions.
No limit is built into these definitions.

\begin{definition}[Separated test functions]
\label{ch01:separated-test}
Set
\[
 \Conf_n=\{(x_1,\ldots,x_n)\in(\RR^4)^n:x_i\ne x_j\text{ for }i\ne j\}.
\]
A separated test function is an element of $C_c^\infty(\Conf_n)$, extended by zero
to $(\RR^4)^n$. It is admissible at volume $p$ if its support is contained
in $((-R_p/2,R_p/2)^4)^n$.
\end{definition}
Such test functions can depend jointly on all positions; they need not factor into
one-point test functions. Compactness of their support gives a positive distance
from the collision diagonals. We write
$\operatorname{sep}_\infty(x)=\min_{i\ne j}|x_i-x_j|_\infty$, with value
$+\infty$ for $n\le1$.

\begin{definition}[Ordered convergence]
\label{ch01:ordered-limit}
A lattice family converges to kernels $S_n(v;x)$ along an increasing
volume-index map $P$ if, for every $n$, every list $v$ and every separated
test function $f$, the following hold. At each admissible fixed volume $p$ there is
a limit
\begin{equation}
 \lim_{N\to\infty}C^{\mathrm{lat}}_{N,p}(v;f)=C_p^{\mathrm{UV}}(v;f),
 \label{ch01:uv-limit}
\end{equation}
and the retained volumes satisfy
\begin{equation}
 \lim_{k\to\infty} C_{P(k)}^{\mathrm{UV}}(v;f)
       =\int_{\Conf_n} f(x)S_n(v;x)\,dx.
 \label{ch01:volume-limit}
\end{equation}
The test function is eventually admissible along $P$.
\end{definition}
The first limit uses every cutoff $N$ in the chosen spacing sequence. The
second may require a subsequence of volumes. The desired construction has
the stronger quantifier order
\begin{enumerate}[label=\textup{(\roman*)}]
\item choose one lattice family;
\item prescribe any increasing sequence of volume indices;
\item extract one refinement that serves every label list and separated test function.
\end{enumerate}
In symbols, for every strictly increasing map $p_{\rm seq}:\NN\to\NN$
there should be a strictly increasing $q_{\rm seq}$ such that
$P=p_{\rm seq}\circ q_{\rm seq}$ works in
Definition~\ref{ch01:ordered-limit}. The lattice family, including its
normalizations, is chosen before $p_{\rm seq}$. Separate subsequences for
different observables would not suffice.

The bare inverse coupling is to satisfy
\[
 \frac{\beta_{N,p}}{\log(1/a_N)}\longrightarrow\kappa_G
 \quad(N\to\infty),\qquad\text{uniformly in }p,
\]
with $\kappa_G$ as in \eqref{ch01:slope}. This is one part of the target,
not a hypothesis that supplies the required continuum bounds. A fixed
parameter $b>0$ specifies the soft constraint in the block transformation;
it is distinct from $\beta$ and from the coefficient $b_0$ in
\eqref{ch01:slope}.

The eventual limit is a family of correlations of gauge-invariant local
fields. It is not presented as a probability measure on continuum
connections. Its decomposition into phases, reconstruction as relativistic
fields, and stress tensor require additional arguments. In particular,
reflection positivity alone does not establish the full reconstruction
theorem. Likewise, a unique vacuum in each reconstructed phase does not
imply a positive mass gap. Neither a mass gap, uniqueness of the
infinite-volume limit, nor independence of the block parameter $b$ is
included in the conclusions described here.

Finally, \eqref{ch01:volume-limit} concerns separated positions. Extending
correlations as distributions across coincident points is a later
renormalization problem. Such an extension is not automatically a limit
of the unsubtracted lattice correlations at collisions. This distinction
will remain important when products of local fields are considered.

\section{How the construction is organized}
\label{ch01:strategy}

The finite Wilson measure is easy to define; its behavior as the lattice
spacing vanishes is the problem. The number of integration variables grows,
and local observables must be rescaled and mixed before their correlations
can have nonzero finite limits. Weak-coupling perturbation theory predicts
some of these changes, but it does not bound the integral over all gauge
configurations. The construction therefore keeps the finite integral exact
while organizing its estimates by scale.

An exact block transformation integrates out information between two
successive lattice spacings and leaves a density on the coarser links.
The transformed density generally has many interactions, even when the
original action is a sum of plaquette terms. The useful question is whether
these new interactions admit a stable description: a prescribed local
action, a small remainder on configurations with small plaquette curvature,
and strongly suppressed contributions where that condition fails.
Here and throughout the book, small curvature means small plaquette
holonomy; it does not mean that every link is close to the identity in an
arbitrary gauge.

These two parts of configuration space require different estimates.
On the small-curvature part, a stationary background and its Gaussian
fluctuations organize the local integral. Taylor coefficients of low order
modify the prescribed action. The other terms must gain enough powers of
the block ratio or the inverse coupling to remain small after summing over
their positions. On the complement, positive action supplies an exponential
penalty. Connected sets of defects are summed with their full combinatorial
multiplicity. This is the role of the polymer expansion: it controls the
sum once the local integral and energy estimates have been proved.
It does not supply those estimates by itself.

The leading change in the coefficient of curvature squared is logarithmic
in the block ratio. Matching this change determines the constant
$\kappa_G$ in \eqref{ch01:slope}. One must then tune a whole finite-cutoff
trajectory so that its coarse endpoint has a prescribed coupling.
The same choice must work with all the estimates that will be required for
local observables. Choosing a different trajectory each time an observable
is introduced would not construct a single quantum field theory.

Local observables bring an additional finite-dimensional problem.
A monomial with $r$ curvature factors and $k$ derivatives has
\emph{scaling dimension} $2r+k$. Gauge-invariant polynomials in curvature and
its derivatives can mix under blocking. A transfer matrix records this mixing at each fixed scaling
dimension. Bounds for this matrix and for the remainder of each inserted
density are needed in addition to bounds for the density without
insertions. Several insertions also require estimates uniform away from
their collision diagonals. Bounds for lattice differences are essential:
bounded correlations alone would not give smooth continuum kernels.

The free Maxwell theory enters as a local comparison at short distances.
Its role is to distinguish field labels, identify possible equation
relations, and normalize certain symmetry identities. Such a comparison
does not assert that all correlations of the interacting theory are
Gaussian. It must be established for the same finite-cutoff family whose
continuum limit is being taken.

After these estimates are available, two separate arguments remove the
cutoffs. At a fixed physical volume, stability of the renormalization
trajectory gives a Cauchy property as the ultraviolet cutoff increases.
Uniform estimates in the physical volume then permit one common
infinite-volume subsequence, as specified in
Definition~\ref{ch01:ordered-limit}. Neither argument can replace the other.

Finally, lattice symmetries and positivity must be promoted to the
continuum statements needed for reconstruction. Coordinate reflection
positivity gives an initial Hilbert-space structure. Bounds on translation
operators, continuous Euclidean symmetries, and control of the correlations
must be established with compatible normalizations. Decomposition into
phases and relativistic reconstruction follow only after these steps.
The trace identity and the expansion of products at short separation
require further estimates of their own. This order explains why the book
keeps the exact finite model, the local analytic estimates, the choice of a
common family, and the continuum consequences distinct.
